% =====================================================================
%  Running Digital Twin - technical report
%  Diego De Nicola, Politecnico di Milano
%
%  Workflow: the text is drafted in Word (draft/report_draft.docx) and
%  then moved here section by section. Equations, figures, tables and
%  references are managed in this file.
%  Compile with pdfLaTeX + BibTeX (Overleaf: default settings).
% =====================================================================
\documentclass[11pt,a4paper]{article}

\usepackage[utf8]{inputenc}
\usepackage[T1]{fontenc}
\usepackage{lmodern}
\usepackage[margin=2.5cm]{geometry}
\usepackage{amsmath,amssymb}
\usepackage{graphicx}
\usepackage{booktabs}
\usepackage{siunitx}
\usepackage[round,authoryear]{natbib}
\usepackage{xcolor}
\usepackage[colorlinks=true,linkcolor=blue!50!black,citecolor=blue!50!black,urlcolor=blue!50!black]{hyperref}

\graphicspath{{figures/}}

% Figure that compiles even if the image file is not there yet
\newcommand{\safefig}[3][0.95\linewidth]{%
  \IfFileExists{figures/#2}%
    {\includegraphics[width=#1]{#2}}%
    {\fbox{\parbox{#1}{\centering\vspace{2cm}\textit{Missing figure: \detokenize{#2}}\vspace{2cm}}}}%
  \caption{#3}}

% Placeholder for text still to be written
\newcommand{\todo}[1]{\textcolor{red!70!black}{\textbf{[TODO:} #1\textbf{]}}}

% Key numbers written by the MATLAB scripts (if available)
\IfFileExists{generated/numbers.tex}{\input{generated/numbers.tex}}{}

\title{Running Digital Twin:\\[4pt]
\large System identification, state estimation and optimal planning\\
of endurance-training adaptation from wearable data}
\author{Diego De Nicola\\
\small MSc Mechanical Engineering (Mechatronics and Robotics), Politecnico di Milano}
\date{\todo{date}}

\begin{document}
\maketitle

% ---------------------------------------------------------------------
\begin{abstract}
\todo{Write LAST, after the race (150--200 words): problem, data, methods
(identification, Kalman filter, MPC), main quantitative results,
pre-registered prediction vs actual half-marathon time.}
\end{abstract}

% ---------------------------------------------------------------------
\section{Introduction}\label{sec:intro}

Every runner who follows a training plan is, without knowing it, controlling a
dynamic system. The input is the training done every day; the output is the
performance on race day. In between there are two quantities that cannot be
measured directly: the fitness built up over weeks and months, and the fatigue
accumulated in the last few days. Hard training raises both. Fitness lasts
long, fatigue goes away quickly, and the performance on a given day depends on
the balance between the two. This is the reason why runners reduce their
training in the last days before a race (the taper): fatigue disappears first,
most of the fitness remains.

Coaches manage this balance with experience. The idea of this project is to
manage it as an engineer: write a model of the athlete, identify its
parameters from data, estimate its hidden states and use it to plan the
training. The athlete is the author, and the data are the ones recorded by his
own sports watch during almost one year of running. The case study is the
preparation of a half marathon on 22~November~2026: the system is used while
the preparation is still running, and the race is its final test.

The project follows the classical chain of a control problem:
\begin{enumerate}
  \item \textbf{Data and input.} A pipeline converts the raw files of the
  watch into a daily training load, computed from heart rate second by second
  (Sections~\ref{sec:data} and~\ref{sec:load}).
  \item \textbf{Output.} Since races are rare, performance is measured on
  ordinary training runs through an aerobic-efficiency index, whose noise is
  quantified and whose dependence on weather and place is checked
  (Section~\ref{sec:perf}).
  \item \textbf{Identification.} The parameters of the fitness--fatigue model
  are estimated from these data, taking advantage of three intentional rest
  periods that act as natural step experiments (Section~\ref{sec:id}).
  \item \textbf{State estimation.} A Kalman filter estimates fitness and
  fatigue day by day as new runs are recorded (Section~\ref{sec:kf}).
  \item \textbf{Planning.} Model predictive control computes the optimal
  training load for the last weeks before the race, and the result is compared
  with the plan actually followed (Section~\ref{sec:mpc}).
  \item \textbf{Test.} The predicted race time is committed to the code
  repository before the race, with a date that cannot be changed, and then
  compared with the real result (Section~\ref{sec:race}).
\end{enumerate}

The last point is what makes the project a real test and not only a fit of
past data: a model adapted after seeing the result can always be made to work,
a prediction fixed in advance cannot. For the same reason, until the race the
model is used only as an advisor, and the training plan is not changed on the
basis of its output.

% ---------------------------------------------------------------------
\section{Background and related work}\label{sec:background}

\subsection{The fitness--fatigue model}
The first systems model of training was proposed by \citet{banister1975} and
described in detail by \citet{calvert1976}. The athlete is treated as a linear
system whose impulse response is the difference of two exponentials: a single
training session produces a positive effect (fitness), large and slow, and a
negative effect (fatigue), smaller in the long run but faster and initially
stronger. Discretised on a daily grid, with training load $w[k]$, the model
reads:
\begin{align}
x_1[k+1] &= e^{-1/\tau_1}\, x_1[k] + w[k] && \text{(fitness)}\label{eq:x1}\\
x_2[k+1] &= e^{-1/\tau_2}\, x_2[k] + w[k] && \text{(fatigue)}\label{eq:x2}\\
p[k]     &= p_0 + k_1 x_1[k] - k_2 x_2[k]  && \text{(performance)}\label{eq:p}
\end{align}
where $\tau_1 > \tau_2$ are the time constants in days, $k_1$ and $k_2$ the
gains and $p_0$ the baseline performance. With the typical values of
$\tau_1 \approx 40$ days and $\tau_2 \approx 7$ days, a single hard session
first makes the athlete worse and only after some days better: this is the
mechanism behind both recovery and taper.

From the engineering point of view, the model is a linear time-invariant
system with one input, two states and one output, and can be written in the
standard state-space form
\begin{equation}
\mathbf{x}[k+1] = \mathbf{A}\,\mathbf{x}[k] + \mathbf{B}\,w[k], \qquad
p[k] = \mathbf{C}\,\mathbf{x}[k] + p_0,
\label{eq:ss}
\end{equation}
with $\mathbf{A} = \mathrm{diag}(e^{-1/\tau_1}, e^{-1/\tau_2})$,
$\mathbf{B} = [1\;\;1]^{\mathsf T}$ and $\mathbf{C} = [k_1\;\;-k_2]$.
This form is the reason why the model fits so naturally with the tools used in
this project: identification, Kalman filtering and model predictive control
all start from a state-space model.

\subsection{Applications and limitations}
\citet{morton1990} applied the model to running performance, and
\citet{fitzclarke1991} used it to find how training sessions should be placed
before a competition, which is the first example of planning based on the
model. Later work focused on its limits. \citet{busso1997} estimated the
parameters recursively, letting them change over time, and \citet{busso2003}
showed that the fatigue response grows with the recent training load, so that
the same session costs more when the athlete is already tired.
\citet{hellard2006} analysed the model on elite athletes and found that its
parameters are estimated with large uncertainty and are strongly correlated
with each other, so that their values must be interpreted with care. Nonlinear
alternatives have also been proposed, such as the PerPot model of
\citet{perl2001} and the model of \citet{turner2017}, which was also used to
optimise the training plan. A clear introduction to the whole topic, with
worked examples, is given by \citet{clarke2013}.

\citet{kolossa2017} combined the fitness--fatigue model with a Kalman filter,
using new performance measurements to correct the state of the model as they
become available. This is the closest work to the state-estimation part of the
present project.

\subsection{Measuring the training load}
The input of the model must be a single number per day. \citet{banister1991}
proposed the training impulse (TRIMP), which weights the duration of the
exercise by the heart-rate reserve with an exponential factor, so that minutes
at high intensity count more than minutes at low intensity. TRIMP is a
measure of \emph{internal} load, i.e.\ of the response of the body, while
distance and pace are measures of \emph{external} load \citep{impellizzeri2019}.
For a model of the response of the athlete, the internal load is the natural
input.

\subsection{Position of this work}
Most studies on the fitness--fatigue model use data collected in laboratory
tests or races, on groups of athletes, and fit the model a posteriori. This
project differs in three respects: it uses only the data of a consumer
sports watch recorded during normal training; it measures performance on
ordinary training runs, so that measurements are frequent; and it uses the
model in real time during a preparation, with a prediction fixed before the
race. The control tools themselves are standard and are taken from the
textbooks of \citet{ljung1999} for identification, \citet{simon2006} for
state estimation and \citet{rawlings2017} for model predictive control.


% ---------------------------------------------------------------------
\section{Data and instrumentation}\label{sec:data}

\subsection{Devices and sensors}
All data were recorded by the author during his own training with two Garmin
smartwatches: a Vivoactive~4S until 2~November~2025 and an Instinct~2~Solar from
3~November~2025 onwards. Heart rate (HR) was always measured by the optical
sensor of the watch at the wrist; no chest strap was used. Since the device was
replaced only two days after the start of the analysed period, the whole
dataset comes in practice from a single device, and no artificial jumps due to a
change of sensor are expected. All processing was implemented in MATLAB.

\subsection{Data processing}
The dataset is the complete export of the author's Garmin account, which
contains 20{,}413 FIT files. Most of them are daily wellness records
(activity monitoring, sleep, heart-rate variability). Activity files were
identified by reading only the header of each file (the \texttt{file\_id}
message), which reduced the files to be converted to 565 activities. These were
converted to text with the FIT CSV Tool \citep{garminfit} and parsed into
per-second records of HR, speed, distance, altitude and GPS position. After
selecting the running activities, 114 runs between 1~November~2025 and
18~October~2026 were retained.

Three data-quality issues were handled. Activities appearing twice in the
data (same start time, different file names) were removed. Activity files
without per-second records, such as manually entered activities, were
discarded. Finally, two runs were recorded on a treadmill: without GPS, their
speed is estimated from the wrist accelerometer and is not comparable with that
of outdoor runs. These two runs are kept in the training load, since their HR
is measured as in any other run, but are excluded from the performance
measurement (Section~\ref{sec:perf}). No run had less than \SI{90}{\percent} of
usable HR samples.

The pipeline was validated on four reference sessions (an easy run, an interval
session, a long run and a \SI{10}{km} race): duration, distance and average HR
computed from the per-second records match the activity summaries of Garmin
Connect (Table~\ref{tab:loadcompare}).

\subsection{Data windows}
Fitness and fatigue are not measurable, and their value at the beginning of
the data is unknown. The series therefore starts on 1~November~2025, four
months before the first data used for identification (Table~\ref{tab:windows}).
Since the effect of the unknown initial state decays as $e^{-t/\tau_1}$, with
the literature value $\tau_1 \approx \SI{42}{days}$ it is reduced to about
\SI{6}{\percent} after four months. Running in this period was not structured,
but it is still training load and is used as input. From 17~August~2026 the
author followed a structured plan towards a \SI{10}{km} race (27~September,
40:59) and then towards a half marathon on 22~November~2026. During this last
period the model is used online; it acts only as an advisor, and the training
plan was not modified on the basis of its output.

\begin{table}[ht]
\centering
\caption{Data windows and their role.}\label{tab:windows}
\begin{tabular}{@{}ll@{}}
\toprule
Window & Role \\
\midrule
1 Nov 2025 -- 28 Feb 2026 & burn-in of the model states \\
1 Mar 2026 -- 16 Aug 2026 & identification (incl.\ summer break 1--16 Aug) \\
17 Aug 2026 -- 19 Oct 2026 & identification and validation (structured plan) \\
20 Oct 2026 -- 22 Nov 2026 & online use; the half marathon is the final test \\
\bottomrule
\end{tabular}
\end{table}

\subsection{Rest periods and places}
The data contain three rest periods: late December~2025, late February to
early March~2026, and 1--16~August~2026. All three were intentional and not
caused by illness or injury. For the model they are therefore true zeros of the
input, not missing data: during them fitness and fatigue decay freely, which
makes these segments natural step-down experiments and particularly
informative about the two time constants.

From the GPS position at the start of each run, the runs were assigned to two
places: Milan (about \SIrange{120}{150}{m} above sea level, flat; 72~runs) and
L'Aquila (about \SIrange{600}{720}{m} above sea level, hilly; 39~runs,
including the whole summer from 3~July to 12~September~2026). The remaining
runs are one run abroad and the two treadmill runs. The effect of place and
weather on the performance measurement is analysed in Section~\ref{sec:perf}.

% ---------------------------------------------------------------------
\section{Training load}\label{sec:load}

The input of the model, $w[k]$, is the daily training load, computed as the
Banister training impulse \citep{banister1991}:
\begin{equation}
\mathrm{TRIMP} = \sum_i \frac{\Delta t_i}{60}\, \mathrm{HR_r}_{,i}\; a\, e^{\,b\, \mathrm{HR_r}_{,i}},
\qquad
\mathrm{HR_r} = \frac{\mathrm{HR}-\mathrm{HR_{rest}}}{\mathrm{HR_{max}}-\mathrm{HR_{rest}}},
\label{eq:trimp}
\end{equation}
with $a=0.64$ and $b=1.92$ (male coefficients), $\Delta t_i$ in seconds.
The resting HR, $\mathrm{HR_{rest}} = \SI{45}{bpm}$, is the seven-day average
measured by the watch in October~2026; the maximum HR,
$\mathrm{HR_{max}} = \SI{200}{bpm}$, is the highest reliable value recorded,
reached during the \SI{10}{km} race.

Because of the exponential weighting, TRIMP is a convex function of HR.
Computed from the average HR of a session, it would therefore underestimate
sessions in which HR varies widely, such as intervals and races. For this
reason it is integrated sample by sample at \SI{1}{Hz}, skipping samples below
\SI{60}{bpm} (sensor dropouts) and gaps longer than \SI{10}{s} (pauses). The
daily load is the sum of the TRIMP of the runs of that day; rest days have
$w[k]=0$.

\begin{table}[ht]
\centering
\caption{Banister TRIMP and Garmin training load for the four reference
sessions, ordered by intensity.}
\label{tab:loadcompare}
\begin{tabular}{@{}llccccc@{}}
\toprule
Date & Session & Duration & Avg HR & TRIMP & Garmin load & Garmin/TRIMP \\
     &         & [min]    & [bpm]  &       &             &              \\
\midrule
13 Oct 2026 & easy run, \SI{9}{km}      & 53.4 & 145 &  77.8 &  79.8 & 1.03 \\
18 Oct 2026 & long run, \SI{17}{km}     & 90.7 & 151 & 152.7 & 152.9 & 1.00 \\
15 Oct 2026 & intervals, $4\times\SI{1}{km}$ & 53.0 & 149 &  88.2 & 168.5 & 1.91 \\
27 Sep 2026 & \SI{10}{km} race          & 41.0 & 188 & 142.0 & 374.8 & 2.64 \\
\bottomrule
\end{tabular}
\end{table}

Table~\ref{tab:loadcompare} compares TRIMP with the training load computed by
Garmin, which is based on an estimate of the excess post-exercise oxygen
consumption (EPOC). The two measures agree on aerobic runs (ratio close to~1)
but diverge above the threshold, and the ratio grows with intensity, up to 2.6
for the race. The interval session is the clearest case: its average HR is
almost the same as that of the long run (149 versus \SI{151}{bpm}), yet Garmin
assigns it about twice the load per minute, because it accounts for the peaks
in the high-intensity zone. Neither measure is the \emph{true} load: they are
two different hypotheses on how intensity is converted into a training
stimulus. They are compared as model inputs in Section~\ref{sec:id}.

\begin{figure}[ht]
\centering
\safefig{fig01_daily_load.pdf}{Daily training load $w[k]$ (TRIMP) and
exponentially weighted moving averages with the literature time constants
of fitness ($\tau = \SI{42}{days}$) and fatigue ($\tau = \SI{7}{days}$).}
\label{fig:load}
\end{figure}

Figure~\ref{fig:load} shows the daily load together with its exponentially
weighted moving averages, computed with the literature time constants of
fitness and fatigue. The fitness proxy starts from zero in November, as
expected during the burn-in, and is already at about 15 in March. It peaks in
July (about~39), drops during the August break (about~24) and is back to
about~36 in mid-October. In each rest period the fatigue proxy decays within a
few days, while the fitness proxy decays much more slowly. This is the
two-time-constant behaviour of the model applied to the measured input, not
yet a validation of the model: the curves are computed with literature
parameters and are not compared with any measurement of performance. Measuring
performance and identifying the parameters are the subject of
Sections~\ref{sec:perf} and~\ref{sec:id}.


% ---------------------------------------------------------------------
\section{Performance measurement}\label{sec:perf}
% Aerobic efficiency on easy steady runs; grade adjustment (Minetti 2002);
% analysis window 10-45 min; selection criteria and sensitivity analysis;
% exclusion of treadmill runs; noise estimate (2.9 %); weather (Open-Meteo,
% ERA5) and place covariates; no day-to-day heat effect; L'Aquila.
Performance is measured by the aerobic efficiency of easy, steady runs,
i.e.\ the grade-adjusted speed per net heartbeat:
\begin{equation}
\mathrm{EF} = \frac{\overline{v}_{\mathrm{GAP}}\;[\si{m/min}]}{\overline{\mathrm{HR}} - \mathrm{HR_{rest}}},
\qquad
v_{\mathrm{GAP}} = v\, \frac{C(i)}{C(0)},
\label{eq:ef}
\end{equation}
where $C(i)$ is the energy cost of running on a gradient $i$
\citep{minetti2002}. Weather data at the place and time of every run are
taken from the Open-Meteo archive \citep{openmeteo2023} based on ERA5
\citep{era5}.
\todo{Performance measurement text.}

\begin{table}[ht]
\centering
\caption{Sensitivity of the run selection to the steady-share threshold.}
\label{tab:sensitivity}
\begin{tabular}{@{}ccccc@{}}
\toprule
min.\ steady share & selected & L'Aquila & Jul--Sep & $\sigma$ [m/beat] \\
\midrule
0.60 & 57 & 12 & 10 & 0.046 \\
0.70 & 57 & 12 & 10 & 0.046 \\
0.80 & 55 & 10 &  9 & 0.046 \\
\bottomrule
\end{tabular}
\end{table}

\begin{figure}[ht]
\centering
\safefig{fig02_efficiency.pdf}{Aerobic efficiency of easy steady runs over
time (raw and temperature-corrected running medians) and fitness proxy.}
\label{fig:ef}
\end{figure}

\begin{figure}[ht]
\centering
\safefig[0.7\linewidth]{fig03_heat_effect.pdf}{EF residual with respect to
a 9-run running median versus air temperature, by place.}
\label{fig:heat}
\end{figure}

% ---------------------------------------------------------------------
\section{System identification}\label{sec:id}
% Step 3. Nonlinear least squares, initial state, identifiability,
% validation, TRIMP vs Garmin load as input, joint temperature term.
\todo{Step 3.}

% ---------------------------------------------------------------------
\section{State estimation}\label{sec:kf}
% Step 4. State-space form, Kalman filter, measurement noise R estimated
% by maximum likelihood, daily online use.
\todo{Step 4.}

% ---------------------------------------------------------------------
\section{Optimal taper planning}\label{sec:mpc}
% Step 5. MPC formulation, constraints, comparison with the actual plan
% \citep{mujika2003}.
\todo{Step 5.}

% ---------------------------------------------------------------------
\section{Pre-registered prediction and race result}\label{sec:race}
% Prediction committed to the repository before 22 Nov 2026 (commit hash
% and date), actual result, error analysis.
\todo{After the race.}

% ---------------------------------------------------------------------
\section{Discussion and limitations}\label{sec:discussion}
% Single athlete (n = 1); wrist optical HR; identifiability; seasonal
% confounding (heat acclimatisation vs fatigue); selection of runs.
\todo{Discussion.}

% ---------------------------------------------------------------------
\section{Conclusions and future work}\label{sec:conclusions}
% Dynamic HR model per run (EF from all runs); chest-strap validation;
% time-varying parameters.
\todo{Conclusions.}

% ---------------------------------------------------------------------
\section*{Reproducibility}
All code is available at \todo{GitHub URL}. Raw data are not published
because activity files contain GPS tracks.

\bibliographystyle{plainnat}
\bibliography{references}

\end{document}
